\textbf{Phần 1. Trắc nghiệm lựa chọn}
\setcounterpageref{ex}{0}

\begin{ex} %[2D1V03]
Trong một phép cộng có hai số hạng, nếu tăng số hạng thứ nhất thêm $8$ đơn vị và giữ nguyên số hạng thứ hai thì tổng sẽ
\choice
{\True Tăng thêm $8$ đơn vị}
{Giảm đi $8$ đơn vị}
{Không đổi}
{Tăng thêm $16$ đơn vị}
\loigiai{
Khi thêm vào một số hạng bao nhiêu đơn vị thì tổng cũng tăng thêm bấy nhiêu đơn vị. Vậy tổng tăng $8$ đơn vị.
\textbf{Đáp án của câu hỏi là A}
}
\end{ex}

\begin{ex} %[2D1V03]
Trong một phép trừ, số bị trừ là $75$, hiệu là $28$. Số trừ là
\choice
{$57$}
{$48$}
{\True $47$}
{$43$}
\loigiai{
Số trừ = Số bị trừ - Hiệu = $75 - 28 = 47$.
\textbf{Đáp án của câu hỏi là C}
}
\end{ex}

\begin{ex} %[2D1V03]
Biết tổng của hai số là $64$. Nếu một số hạng là $20$ thì số hạng còn lại là
\choice
{$34$}
{\True $44$}
{$54$}
{$42$}
\loigiai{
Số hạng còn lại = Tổng - Số hạng đã biết = $64 - 20 = 44$.
\textbf{Đáp án của câu hỏi là B}
}
\end{ex}

\begin{ex} %[2D1V03]
Trong phép trừ, nếu tăng số bị trừ thêm $5$ đơn vị và tăng số trừ thêm $5$ đơn vị thì hiệu sẽ
\choice
{Tăng thêm $10$ đơn vị}
{Giảm đi $5$ đơn vị}
{Tăng thêm $5$ đơn vị}
{\True Không thay đổi}
\loigiai{
Khi cùng thêm (hoặc bớt) ở cả số bị trừ và số trừ một số đơn vị như nhau thì hiệu không thay đổi.
\textbf{Đáp án của câu hỏi là D}
}
\end{ex}

\begin{ex} %[2D1V03]
Tìm một số biết rằng lấy số đó cộng với $18$ rồi trừ đi $12$ thì được $45$.
\choice
{\True $39$}
{$41$}
{$37$}
{$43$}
\loigiai{
Tính ngược từ cuối: $45 + 12 - 18 = 57 - 18 = 39$.
\textbf{Đáp án của câu hỏi là A}
}
\end{ex}

\begin{ex} %[2D1V03]
Trong một phép trừ, tổng của số bị trừ, số trừ và hiệu bằng $70$. Số bị trừ là
\choice
{$40$}
{$30$}
{\True $35$}
{$45$}
\loigiai{
Vì Số trừ + Hiệu = Số bị trừ nên: Số bị trừ + Số bị trừ = 70. Vậy Số bị trừ là $70 : 2 = 35$.
\textbf{Đáp án của câu hỏi là C}
}
\end{ex}

\begin{ex} %[2D1V03]
An có $32$ viên bi, Bình có $18$ viên bi. Để hai bạn có số bi bằng nhau thì An phải cho Bình bao nhiêu viên bi?
\choice
{$14$ viên bi}
{\True $7$ viên bi}
{$8$ viên bi}
{$6$ viên bi}
\loigiai{
An hơn Bình: $32 - 18 = 14$ viên bi.\nAn phải cho Bình: $14 : 2 = 7$ viên bi.
\textbf{Đáp án của câu hỏi là B}
}
\end{ex}

\begin{ex} %[2D1V03]
Tìm $x$ biết: $x - 24 = 35 + 11$.
\choice
{$60$}
{$68$}
{$72$}
{\True $70$}
\loigiai{
Ta có: $x - 24 = 46 \Rightarrow x = 46 + 24 = 70$.
\textbf{Đáp án của câu hỏi là D}
}
\end{ex}

\begin{ex} %[2D1V03]
Hai số có tổng bằng $58$. Nếu bớt ở số hạng thứ nhất đi $6$ đơn vị và thêm vào số hạng thứ hai $9$ đơn vị thì tổng mới là
\choice
{\True $61$}
{$55$}
{$63$}
{$59$}
\loigiai{
Tổng mới là: $58 - 6 + 9 = 61$.
\textbf{Đáp án của câu hỏi là A}
}
\end{ex}

\begin{ex} %[2D1V03]
Trong phép trừ có hiệu bằng $0$. Khẳng định nào sau đây luôn đúng?
\choice
{Số bị trừ lớn hơn số trừ}
{Số trừ bằng $0$}
{\True Số bị trừ bằng số trừ}
{Số bị trừ bằng $0$}
\loigiai{
Hiệu bằng 0 khi và chỉ khi số bị trừ bằng số trừ.
\textbf{Đáp án của câu hỏi là C}
}
\end{ex}

\begin{ex} %[2D1V03]
Mai thực hiện phép tính: $45 - a = 18$. Số $a$ đóng vai trò là gì và có giá trị là bao nhiêu?
\choice
{Số bị trừ, giá trị là $63$}
{\True Số trừ, giá trị là $27$}
{Hiệu, giá trị là $27$}
{Số hạng, giá trị là $18$}
\loigiai{
$a$ là số trừ trong phép trừ. Giá trị của $a = 45 - 18 = 27$.
\textbf{Đáp án của câu hỏi là B}
}
\end{ex}

\begin{ex} %[2D1V03]
Một sợi dây dài $85\text{ cm}$. Người ta cắt đi một đoạn thì còn lại đoạn dây dài $38\text{ cm}$. Độ dài đoạn dây bị cắt đi là
\choice
{$57\text{ cm}$}
{$43\text{ cm}$}
{$53\text{ cm}$}
{\True $47\text{ cm}$}
\loigiai{
Đoạn dây bị cắt đi là: $85 - 38 = 47\text{ (cm)}$.
\textbf{Đáp án của câu hỏi là D}
}
\end{ex}

\textbf{Phần 2. Trắc nghiệm đúng sai}
\setcounterpageref{ex}{0}

\begin{ex} %[2D1V03]
Xét tính đúng sai của các nhận định về phép cộng và phép trừ:
\choiceTF
{\True Trong phép cộng, khi đổi chỗ các số hạng thì tổng không thay đổi.}
{Trong phép trừ, số bị trừ luôn luôn phải lớn hơn số trừ.}
{\True Số bị trừ bằng tổng của số trừ và hiệu.}
{Nếu tổng bằng 0 thì cả hai số hạng đều phải bằng 1.}
\loigiai{
\begin{itemize}
    \item Ý a) Đúng: Đây là tính chất giao hoán của phép cộng.
    \item Ý b) Sai: Số bị trừ có thể bằng số trừ (khi đó hiệu bằng 0).
    \item Ý c) Đúng: Công thức cơ bản: Số bị trừ = Số trừ + Hiệu.
    \item Ý d) Sai: Trong tập số tự nhiên, nếu tổng bằng 0 thì cả hai số hạng đều phải bằng 0.
\end{itemize}
\begin{itemize}
    \item a) \textbf{Đúng}.
    \item b) \textbf{Sai}.
    \item c) \textbf{Đúng}.
    \item d) \textbf{Sai}.
\end{itemize}
}
\end{ex}

\begin{ex} %[2D1V03]
Cho phép tính $A - B = C$ (với $A, B, C$ là các số tự nhiên khác 0):
\choiceTF
{\True Nếu giữ nguyên $A$ và tăng $B$ thêm 3 đơn vị thì $C$ giảm đi 3 đơn vị.}
{\True Nếu tăng $A$ thêm 4 đơn vị và giữ nguyên $B$ thì $C$ tăng thêm 4 đơn vị.}
{Nếu $A$ tăng 5 đơn vị và $B$ giảm 5 đơn vị thì $C$ không đổi.}
{\True Tổng $A + B + C$ luôn là một số chẵn.}
\loigiai{
\begin{itemize}
    \item Ý a) Đúng: Khi số trừ tăng thì hiệu giảm tương ứng.
    \item Ý b) Đúng: Khi số bị trừ tăng thì hiệu tăng tương ứng.
    \item Ý c) Sai: Hiệu mới sẽ tăng thêm $5 + 5 = 10$ đơn vị.
    \item Ý d) Đúng: Vì $B + C = A$ nên $A + B + C = A + A = 2 \times A$, luôn là số chẵn.
\end{itemize}
\begin{itemize}
    \item a) \textbf{Đúng}.
    \item b) \textbf{Đúng}.
    \item c) \textbf{Sai}.
    \item d) \textbf{Đúng}.
\end{itemize}
}
\end{ex}

\textbf{Phần 3. Trả lời ngắn}
\setcounterpageref{ex}{0}

\begin{ex} %[2D1V03]
Trong một phép trừ, số bị trừ hơn số trừ 26 đơn vị. Hỏi hiệu của phép trừ đó bằng bao nhiêu?
\shortans{26}
\loigiai{
\textbf{Đáp số:} $26$\\
Hiệu chính là phần số bị trừ hơn số trừ, do đó hiệu bằng 26.
}
\end{ex}

\begin{ex} %[2D1V03]
Tìm một số biết rằng lấy 78 trừ đi số đó rồi cộng thêm 15 thì được 43.
\shortans{50}
\loigiai{
\textbf{Đáp số:} $50$\\
Ta có: $78 - x = 43 - 15 = 28 \Rightarrow x = 78 - 28 = 50$.
}
\end{ex}

\begin{ex} %[2D1V03]
Tổng của hai số là 75. Nếu số hạng thứ nhất tăng thêm 12 đơn vị và số hạng thứ hai giảm đi 7 đơn vị thì tổng mới bằng bao nhiêu?
\shortans{80}
\loigiai{
\textbf{Đáp số:} $80$\\
Tổng mới: $75 + 12 - 7 = 80$.
}
\end{ex}

\begin{ex} %[2D1V03]
Minh có 28 cái kẹo, Lan có 16 cái kẹo. Mẹ cho thêm hai bạn một số kẹo như nhau thì tổng số kẹo của hai bạn là 50 cái. Hỏi mẹ đã cho mỗi bạn bao nhiêu cái kẹo?
\shortans{3}
\loigiai{
\textbf{Đáp số:} $3$\\
Tổng số kẹo ban đầu là $28 + 16 = 44$. Mẹ cho cả hai bạn $50 - 44 = 6$ cái kẹo. Mỗi bạn được $6 : 2 = 3$ cái.
}
\end{ex}

\begin{ex} %[2D1V03]
Tìm số trừ trong một phép trừ biết rằng số bị trừ là 90 và hiệu là số chẵn nhỏ nhất có hai chữ số.
\shortans{80}
\loigiai{
\textbf{Đáp số:} $80$\\
Số chẵn nhỏ nhất có hai chữ số là 10. Số trừ là $90 - 10 = 80$.
}
\end{ex}

\begin{ex} %[2D1V03]
Hai túi gạo có tất cả 64 kg gạo. Nếu chuyển 6 kg gạo từ túi thứ nhất sang túi thứ hai thì hai túi có số gạo bằng nhau. Hỏi lúc đầu túi thứ nhất có bao nhiêu ki-lô-gam gạo?
\shortans{38}
\loigiai{
\textbf{Đáp số:} $38$\\
Túi 1 hơn túi 2 là $6 \times 2 = 12\text{ kg}$. Lúc đầu túi 1 có: $(64 + 12) : 2 = 38\text{ kg}$.
}
\end{ex}
